\documentclass[fleqn]{article}
\oddsidemargin 0.0in
\textwidth 6.0in
\thispagestyle{empty}
\usepackage{import}
\usepackage{amsmath}
\usepackage{graphicx}
\usepackage{flexisym}
\usepackage{amssymb}
\usepackage{bigints} 
\usepackage[english]{babel}
\usepackage[utf8x]{inputenc}
\usepackage{float}
\usepackage[colorinlistoftodos]{todonotes}

\definecolor{hwColor}{HTML}{AD53BA}

\begin{document}

  \begin{titlepage}

    \newcommand{\HRule}{\rule{\linewidth}{0.5mm}}

    \center



    \textsc{\LARGE Arizona State University}\\[1.5cm]

    \textsc{\LARGE Mathematical Methods For Physics II }\\[1.5cm] 


    \begin{figure}
      \includegraphics[width=\linewidth]{asu.png}
    \end{figure}


    \HRule \\[0.4cm]
    { \huge \bfseries Homework Eight}\\[0.4cm] 
    \HRule \\[1.5cm]

    \textbf{Behnam Amiri}

    \bigbreak

    \textbf{Prof: Cecilia Lunardini}

    \bigbreak


    \textbf{{\large \today}\\[2cm]}

    \vfill

  \end{titlepage}

  \textbf{Part A}
  \begin{enumerate}
    \item Fill the missing steps in sec. 18.5 of the textbook. Starting from eq. (18.74) go as far as proving eqs. (18.77) and (18.78). 
    
      \textcolor{hwColor}{
        Let's start off with the Bessel's differential equation. (first kind) \\
        \\
        $
          \dfrac{d^2 y}{dx^2}+\dfrac{1}{x} \dfrac{dy}{dx}+\left(1-\dfrac{v^2}{x^2}\right)y=0
        $ 
        \\
        \\
        From the textbook (page 603) "By inspection, $x=0$ is a regular sigular point;...", hence, \\
        \\
        \\
        $
          y=x^{\sigma} \sum\limits_{n=0}^{\infty} a_n x^n=\sum\limits_{n=0}^{\infty} a_nx^{n+\sigma} \\
          \\
          \dfrac{dy}{dx}=\dfrac{d}{dx}\left[\sum\limits_{n=0}^{\infty} a_nx^{n+\sigma}\right]=\sum\limits_{n=0}^{\infty} a_n (n+\sigma) x^{n+\sigma-1} \\
          \\
          \\
          \\
          \dfrac{d^2y}{dx^2}=\dfrac{d}{dx}\left[\sum\limits_{n=0}^{\infty} a_n (n+\sigma) x^{n+\sigma-1}\right]=\sum\limits_{n=0}^{\infty} a_n (n^2+\sigma^2+2n\sigma-n-\sigma) x^{n+\sigma-2} \\
          \\
        $ 
        Now it's time to plugin what we found into the first kind Bessel's differential equation. \\
        \\
        $
          \dfrac{d^2 y}{dx^2}+\dfrac{1}{x} \dfrac{dy}{dx}+\left(1-\dfrac{v^2}{x^2}\right)y=0 \\
          \\
          \\
          \left[\sum\limits_{n=0}^{\infty} a_n (n^2+\sigma^2+2n\sigma-n-\sigma) x^{n+\sigma-2}\right]
          +\dfrac{1}{x} \left[\sum\limits_{n=0}^{\infty} a_n (n+\sigma) x^{n+\sigma-1}\right] \\
          +\left(1-\dfrac{v^2}{x^2}\right) \left[\sum\limits_{n=0}^{\infty} a_nx^{n+\sigma} \right]=0 \\
          \\
          \\
          \sum\limits_{n=0}^{\infty} \left[(n+\sigma)-v^2+(n+\sigma-1)(n+\sigma)\right] a_n x^{n+\sigma-2}+\sum\limits_{n=0}^{\infty} a_n x^{n+\sigma}=0 \\
          \\
          \\
          \left[\sum\limits_{n=0}^{\infty} \left[(n+\sigma)-v^2+(n+\sigma-1)(n+\sigma)\right] a_n x^{n+\sigma-2}+\sum\limits_{n=0}^{\infty} a_n x^{n+\sigma}\right] \times x^{2-\sigma}=0 \times x^{2-\sigma} \\
          \\
          \\
          \sum\limits_{n=0}^{\infty}a_n x^{n+2}+\sum\limits_{n=0}^{\infty} \left[(n+\sigma)^2-v^2\right]a_n x^n=0
          \\
          \\
        $
        For the coefficient of x when $n=0$ and $n=1$ we have: \textbf{(The Frobenius solution)} \\ \\
        \emph{Side note: a good way to think about it is that if $a_0$ was 0, the recursion relationship 
        would say that all terms are zero because $a_0$ will be subbed into every other term(the evens at least)
        since we get a contradiction it implies that $a_0 \neq 0 $ so the other portion must} \\ \\
        $
          \begin{cases}
            n=0 ~ \rightarrow ~ \left[(0+\sigma)^2-v^2\right]a_0=\left[\sigma^2-v^2\right]a_0=0 \Rightarrow \sigma=\pm v \\
            \\
            n=1 ~ \rightarrow ~ \left[(1+\sigma)^2-v^2\right]a_1 \Rightarrow \left[(1+\sigma)^2-v^2\right]a_1b \\
          \end{cases} \\
          \\
        $
        For power of 2 and higher, we have this coefficient: \\
        \\
        $
          \left[(n+\sigma)^2-v^2\right]a_n+a_{n-2}=0
        $
        \\
        \\
        A for $x^2$ and higher powers; \\
        \\
        $
          \sum\limits_{n=0}^{\infty}a_n \left[(n+\sigma)^2-v^2\right]+\sum\limits_{n=0}^{\infty} a_{n-2} x^n=0 \\
        $
        \\ 
        \\
        For sigma we can find the recurrence relations as,  \\
        \\
        $
          \begin{cases}
            \left[(1+\sigma)^2-v^2\right]a_1=0 \Rightarrow \left[(1\pm v)^2-v^2\right]a_1=\left[1\pm+2v+v^2-v^2\right]a_1=0 \Rightarrow ~ (1\pm 2v)a_1=0 \\
            \\
            \left[(n+\sigma)^2-v^2\right]a_n+a_{n-2}=0 \Rightarrow \left[n^2 \pm 2nv+v^2-v^2\right]a_n+a_{n-2}=0 \\
          \end{cases} \\
          \\ \\
          \therefore ~~ n(n\pm 2v)a_n+a_{n-2}=0 
        $
         for $n \geq 2 \in \mathcal{Z}$
      }

    \item Continue filling in the steps starting with eq. (18.78) and go as far as proving eq. (18.79) (upper line only). 
    
    
    \item Prove the result at the bottom of the page 606 of the textbook, i.e., that for $\nu$ integer, $J_{-\nu }(x)=(-1)^\nu J_{\nu }(x)$. \\
    (Hint: Start with the series expression for $J_{-\nu }(x)$, which is given by Eq. 18.79 of the textbook where you have to replace $\nu \rightarrow -\nu$. Note that if $k$ is a non-negative integer, then $(-1)^k \Gamma(-k) = \infty$, see fig. 18.9 of the textbook. Eliminate those terms from the sum that equal zero because of this, and then change the summation variable.)
    
  \end{enumerate}


  \textbf{Part B}
  \begin{enumerate}
    \item For a system in spherical polar coordinates, the radial part of Laplace's equation is
    $$r^2 \dfrac{d^2 R}{dr^2}+2r \dfrac{dR}{dr}+\left[k^2r^2-n(n+1)\right]R=0$$
    Using the method shown in class (class 16, example segment on transforming a second order differential equation into Bessel's equation), find it solution in terms of two linearly independent Bessel functions. (Note: write your solution in terms of one Bessel function of the first kind and one Bessel function of the second kind). 

      \textcolor{hwColor}{
        Let's start off by talking a bit about Bessel functions. They are solutions $y(x)$ of
        Bessel's differential equation. \\
        $$x^2 \dfrac{d^2 y}{dx^2}+x \dfrac{dy}{dx}+\left(x^2-\alpha^2\right)y=0$$ \\
        \\
        $$y(x)=W ~ J_{\alpha}+V ~ Y_{\alpha}$$
        \\
        Even though $\alpha$ and $-\alpha$ produce the same differential equation, it is conventional to 
        define different Bessel functions for these two values in such a way that the Bessel functions 
        are mostly smooth functions of $\alpha$. \\
        \\
        The most important cases are when $\alpha$ is an \textbf{integer} or \textbf{half-integer}. Bessel functions for integer 
        $\alpha$ are also known as cylinder functions or the cylindrical harmonics because they appear in the solution 
        to Laplace's equation in cylindrical coordinates. Spherical Bessel functions with half-integer $\alpha$ are obtained 
        when the Helmholtz equation is solved in spherical coordinates. \\
        \\
        \textbf{Bessel functions of the first kind:} $J_{\alpha}$ \\
        Bessel functions of the first kind, denoted as $J_{\alpha}(x)$, are solutions of Bessel's differential equation. 
        For integer or positive $\alpha$, Bessel functions of the first kind are finite at the origin (x = 0); while for negative 
        non-integer $\alpha$, Bessel functions of the first kind diverge as x approaches zero.
        \\
        \\
        \\
        \textbf{Bessel functions of the second kind:} $Y_{\alpha}$ \\
        The Bessel functions of the second kind, denoted by $Y_{\alpha}(x)$, are solutions of the Bessel differential equation
        that have a singularity at the origin (x = 0) and are multivalued. \\ \\
      }


      \textcolor{hwColor}{
        For a Bessel's equation like this $u^2 \dfrac{d^2 f}{du^2}+u\dfrac{df}{du}+\left(u^2-m^2\right)f=0$, we have: \\
        \\
        $
          u=Kx^{\beta}, ~~~~ y(x)=x^{\alpha} f(u) \\ \\
          \begin{cases}
            \dfrac{dy}{dx}=\dfrac{d}{dx}\left(x^{\alpha f}\right)=\alpha x^{\alpha-1}+x^{\alpha} \dfrac{df}{du} \dfrac{du}{dx} \\
            \\
            \dfrac{du}{dx}=\dfrac{d}{dx}\left(C x^{\beta}\right)=\beta Cx^{\beta -1}
          \end{cases} \longrightarrow
          \dfrac{dy}{dx}=x^{\alpha-1} \left[\alpha f+\beta C x^{\alpha} \dfrac{df}{du}\right] \\
          \\
          \\
          \dfrac{d^2 f}{du^2} \times \dfrac{du}{dx}=\dfrac{d}{dx} \left(\dfrac{df}{du}\right)=\dfrac{d}{du} \left(\dfrac{df}{du}\right) \times \dfrac{du}{dx}
          \Longrightarrow \dfrac{1}{\beta C x^{\beta -1}} \dfrac{d}{dx} \left(\dfrac{df}{du}\right) \\
          \\
          \dfrac{d^2 f}{du^2}=\dfrac{1}{\left(\beta C\right)^2 x^{2 \beta+\alpha}} \left(\alpha(\alpha+\beta)y+(1-2\alpha-\beta)x\dfrac{dy}{dx}+x^2\dfrac{d^2 y}{dx^2}\right) \\
          \\
          \\
          \dfrac{1}{\left(\beta C\right)^2 x^{2 \beta+\alpha}} \left(\alpha(\alpha+\beta)y+(1-2\alpha-\beta)x\dfrac{dy}{dx}+x^2\dfrac{d^2 y}{dx^2}\right)
          +\dfrac{C x^{\beta}}{\beta C x^{\alpha+\beta}} \left[x \dfrac{dy}{dx}-\alpha y\right]
          +\left[C^2 x^{2\beta}-m^2\right]\dfrac{y}{x^{\alpha}}=0 \\ \\
        $
        This reduces to: \\
        \\
        $
          x^2 \dfrac{d^2 y}{dx^2}+\left(1-2 \alpha\right)x\dfrac{dy}{dx}+\left[\alpha^2+\beta^2\left(C^2 x^{2\beta}-m^2\right)\right]y=0
        $ \\
        \\
        This equation gives us three parameters $\left(\alpha, ~~ \beta, ~~ and ~~ C\right)$ that we can use to transform our second-order differential equations.
        the goal is to transform the given differntial equation into a Bessel's equation form. Hence, \\
        \\
        \\
        $
          \begin{cases}
            x^2 \dfrac{d^2 y}{dx^2}+\left(1-2 \alpha\right)x\dfrac{dy}{dx}+\left[\alpha^2+\beta^2\left(C^2 x^{2\beta}-m^2\right)\right]y=0 \\
            \\
            r^2 \dfrac{d^2 R}{dr^2}+2r \dfrac{dR}{dr}+\left[k^2r^2-n(n+1)\right]R=0
          \end{cases} \Rightarrow 1-2 \alpha=2 \Rightarrow \alpha=-\dfrac{1}{2}
        $ \\
        \\ 
        \\
        Now, therefore: \\
        \\
        $
          r^2 \dfrac{d^2 R}{dr^2}+2r \dfrac{dR}{dr}+\left[\dfrac{1}{4}+\beta^2\left(C^2 x^{2\beta}-m^2\right)\right]R=0
        $ \\
        \\
        By equating, $\left[\dfrac{1}{4}+\beta^2\left(C^2 x^{2\beta}-m^2\right)\right]R$ to the given P.D.E, we get: \\
        \\
        $
          \left[\dfrac{1}{4}+\beta^2\left(C^2 x^{2\beta}-m^2\right)\right]R=\left[k^2r^2-n(n+1)\right]R \Longrightarrow 
          \dfrac{1}{4}+\beta^2\left(C^2 x^{2\beta}-m^2\right)=k^2r^2-n(n+1) \\
          \\
          \\
          \therefore ~~ \dfrac{1}{4}+\beta^2 C^2 x^{2\beta}-\beta^2 m^2=k^2r^2-n(n+1) \\
          \\
        $
        Asssuming $\beta=1$ we have $\dfrac{1}{4}+C^2 x^2-m^2=k^2r^2-n(n+1)$.\\ \\
        $
          m^2-\dfrac{1}{4}=n(n+1) \\
          \\
          \therefore ~~ m=\pm \left(n+\dfrac{1}{2}\right) ~~~.
        $ (Note: $k^2 r^2=C^2 x^2$) \\
        \\
        We are almost done since we found all the three parameters: \\
        \\
        $
          \begin{cases}
            \alpha=-\dfrac{1}{2} \\
            \\
            \beta=1 \\
            \\
            C=\pm k \\
            \\
            m=\left(n+\dfrac{1}{2}\right)
          \end{cases}
        $
        \\
        \\
        \\
        As we mentioned earlier the solution for the given differential equation is $ y(x)=x^{\alpha} f(u)$, therefore: \\
        \\
        $
          y(x)=y(x)=x^{\alpha} f(u)=W x^{\alpha} J_m \left(K x^{\beta}\right)+V x^{\alpha} Y_m \left(Kx^{\beta}\right) \\
          \\
          =W x^{-\dfrac{1}{2}} J_{m} \left(K x\right)+V x^{-\dfrac{1}{2}} Y_{m} \left(Kx\right) \\
          \\
          \\
          \\
          \therefore ~~~~ y(x)=W x^{-\dfrac{1}{2}} J_{n+\dfrac{1}{2}} \left(K x\right)+V x^{-\dfrac{1}{2}} Y_{n+\dfrac{1}{2}} \left(Kx\right) \\
          \\
        $
        $W$ and $V$ can be determined by boundry conditions.
      }

    \item Using the series expression of the Bessel function $J_{\nu}$, prove eq. (18.93) of the textbook. 

      \textcolor{hwColor}{
        $
          \dfrac{d}{dx} \left(x^{-v} \mathcal{J}_{v}(x)\right)=-x^{-v} \mathcal{J}_{v+1}(x) \\
          \\
          \\
          =\dfrac{d}{dx}\left[x^{-v}\sum\limits_{n=0}^{\infty} \dfrac{(-1)^n}{n!\Gamma\left(1+v+n\right) }(\dfrac{x}{2})^{2n+v}\right]=\sum\limits_{n=0}^{\infty} \dfrac{(-1)^n}{n!\Gamma\left(1+v+n\right)} \dfrac{d}{dx}\left(\dfrac{x^{2n}}{2^{2n+v}}\right) \\
          \\
          \\
          =\sum\limits_{n=0}^{\infty} \dfrac{(-1)^n}{n!\Gamma\left(1+v+n\right)} \dfrac{1}{2^{2n+v}}\dfrac{d}{dx}\left(x^{2n}\right)=\sum\limits_{n=0}^{\infty} \dfrac{(-1)^n}{n!\Gamma\left(1+v+n\right)} \dfrac{2n}{2^{2n+v}}x^{2n-1} \\
          \\
          \\
          =\sum\limits_{n=0}^{\infty} \dfrac{(-1)^n}{n!\Gamma\left(1+v+n\right)} \dfrac{n}{2^{2n+v-1}}x^{2n-1}=\sum\limits_{n=0}^{\infty} \dfrac{(-1)^n}{(n-1)!\Gamma\left(1+v+n\right)} \dfrac{1}{2^{2n+v-1}}x^{2n-1} \\
          \\
          \\
          =x^{-v}\sum\limits_{n=0}^{\infty} \dfrac{(-1)^n}{(n-1)!\Gamma\left(1+v+n\right)} \left(\dfrac{x}{2}\right)^{2n+v-1} \\
          \\
          \\
          =x^{-v}\sum\limits_{n=0}^{\infty} \dfrac{(-1)^n}{((u+1)-1)!\Gamma\left(1+v+(u+1)\right)} \left(\dfrac{x}{2}\right)^{2(u+1)+v-1}=-x^{-v}  \mathcal{J}_{v+1}(x) \\
          \\
          \\
          \\
          \\
          \therefore ~~~  \dfrac{d}{dx} \left(x^{-v} \mathcal{J}_{v}(x)\right)=-x^{-v} \mathcal{J}_{v+1}(x)
        $
      }


    \item Using a graphing tool of your choice (e.g., Wolfram alpha, where all you need is a browser), plot $J_4$ and $J_5$, and verify that the $n-$th root of $J_5$ is always larger than the $n-$th root of $J_4$ (for the same $n$; do not count $x=0$ as a root).  In other words, that $\zeta_{4n}<\zeta_{5n}$ for all $n$. (include the graph in your submission; it will suffice to check the result for $n=1,2,3,4$). 


  \end{enumerate}

\end{document}
